\chapter{EXPECTED RESULTS}

\section{Expected Results of the Research}

A comparative analysis between the RT-LIVO2SM approach and existing methods is expected to demonstrate that RT-LIVO2SM can deliver advantageous outcomes in real-time object detection scenarios when compared to existing methods.

\section{Preliminary Results}

\sloppy
Preliminary research has uncovered that Waymo Open Dataset preferrable for benchmarking the performance of the RT-LIVO2SM system, due to the large number of previous implementations that have published their performance results using this dataset. Communication with the Robotics team at ITS, specifically Team Barunastra ITS, has been done regarding the potential provisioning of the NVIDIA Jetson AGX Orin Developer Kit for performance testing. Additionally, practical unlimited access to the computational capabilities of the NVIDIA AD103-400 for model training during the research period has been secured, ensuring minimal disruptions in model training.

As a footnote, the researcher already has experience in performing \emph{debugging} and \emph{tweaking} on the DLIO architecture \parencite{dlio}. With approximately two years of active experience in dataset collection and processing for \emph{Computer Vision} via the Roboflow platform—gained during their time as a \emph{Perceptions Programmer} for the Barunastra ITS robotics team—the researcher has also experimented with the YOLOv5, YOLOv7, YOLOv8, YOLOv9, YOLOv10, and YOLOv11 architectures. This background gives the researcher confidence in augmenting the capabilities of both DLIO and YOLO architectures by integrating them into the proposed RT-LIVO2SM architecture.
